\documentclass[../../main.tex]{subfiles}
\begin{document}

% 年度の大きな表示
\begin{center}
  {\fontsize{48pt}{58pt}\selectfont \textbf{2008}\normalsize\textbf{年度}}
\end{center}

\vspace{10mm}

% 本文


\vspace{15mm}

% 出題テーマと難易度の表
\noindent\textbf{出題テーマと難易度}

\vspace{5mm}

\begin{center}
  \shadowbox{%
    \begin{tabular}{|c|p{8cm}|c|}
      \hline
      \textbf{問題番号} & \textbf{テーマ} & \textbf{難易度} \\
      \hline
      第1問 & 曲線の接する条件と面積の極限（べき関数と対数関数） & \\
      \hline
      第2問 & ガウス記号を含む極限の場合分け & \\
      \hline
      第3問 & 確率の不等式証明（コーシー・シュワルツ） & \\
      \hline
      第4問 & 正三角形の頂点の軌跡（楕円） & \\
      \hline
    \end{tabular}%
  }
\end{center}

\vspace{15mm}

% 解答
\noindent\textbf{解答}

\vspace{5mm}

\begin{center}
  \shadowbox{%
    \renewcommand{\arraystretch}{1.6}
    \begin{tabular}{|c|c|p{8cm}|}
      \hline
      \textbf{問題番号} & \textbf{小問} & \textbf{答え} \\
      \hline
      \multirow{2}{*}{第1問} & (1) & $s=\left(\dfrac{1}{a}\right)^{1/a}$，$t=\dfrac{1}{a}$，$b=(ae)^{1/a}$ \\
      \cline{2-3}
                             & (2) & $\dfrac{a}{a+1}\left(\dfrac{1}{a}\right)^{1/a}$ \\
      \hline
      第2問 & - & $a\neq b$のとき $c$の最大値$1$，極限値$\dfrac{b-a}{ab}$；$a=b$のとき $c$の最大値$2$，極限値$\dfrac{10}{a^2}$ \\
      \hline
      \multirow{2}{*}{第3問} & (1) & 証明問題 \\
      \cline{2-3}
                             & (2) & 証明問題 \\
      \hline
      \multirow{2}{*}{第4問} & (1) & $R\left(\dfrac{1}{2\tan\alpha}+\dfrac{\sqrt{3}}{2},\,0\right)$ \\
      \cline{2-3}
                             & (2) & 証明問題 \\
      \hline
    \end{tabular}%
    \renewcommand{\arraystretch}{1.0}
  }
\end{center}

\vspace{15mm}

% 解答時間
\noindent\textbf{試験形態}

\vspace{5mm}

\begin{center}
  \shadowbox{%
    \begin{tabular}{p{8cm}}
      （要確認）
    \end{tabular}%
  }
\end{center}

\end{document}
